\documentclass[aspectratio=169,11pt]{beamer}
\usepackage[T1]{fontenc}
\usepackage[utf8]{inputenc}
\usepackage[french]{babel}
\usepackage{lmodern}
\usepackage{booktabs}
\usepackage{tabularx}

% Réunion d'avancement ; durée à ajuster selon les nouveaux papiers et la démonstration.
% Sources locales : slides_6_oct.tex ; etat_de_l_art_attaques_ml_nids.tex ;
% ../docs/attaques_concretes_ML_NIDS/selection_attaques.md ; ../code/README.md.
% Les fonctions du prototype sont documentées et présentes dans le code ;
% cette présentation ne prétend pas à une nouvelle validation en réseau actif.
\definecolor{projectblue}{HTML}{17365D}
\definecolor{projectteal}{HTML}{087E8B}
\usetheme{Madrid}
\setbeamercolor{structure}{fg=projectblue}
\setbeamercolor{frametitle}{bg=projectblue,fg=white}
\setbeamercolor{block title}{bg=projectteal,fg=white}
\setbeamercolor{block body}{bg=projectteal!7}
\setbeamertemplate{navigation symbols}{}
\setbeamertemplate{itemize items}[circle]
\setbeamertemplate{footline}{%
  \leavevmode\hbox{%
    \begin{beamercolorbox}[wd=.84\paperwidth,ht=2.7ex,dp=1.2ex,leftskip=1em]{author in head/foot}%
      VAP SSR 2026--2027 \hfill Point d'avancement -- 27 octobre 2026
    \end{beamercolorbox}%
    \begin{beamercolorbox}[wd=.16\paperwidth,ht=2.7ex,dp=1.2ex,center]{date in head/foot}%
      \insertframenumber/\inserttotalframenumber
    \end{beamercolorbox}}%
}
\newcommand{\espacement}{\setlength{\itemsep}{.8em}}
\title[Démonstrateur ML-NIDS]{Démonstrateur d'attaques contre un ML-NIDS}
\subtitle{État de l'art et premier prototype réseau}
\author{Projet VAP SSR 2026--2027}
\institute{Réunion avec Grégory Blanc}
\date{27 octobre 2026}

\begin{document}

\begin{frame}
  \titlepage
  \note{Présenter les deux avancées : la synthèse bibliographique et le site de construction réseau. L'intégration ML constitue le prochain travail.}
\end{frame}

\begin{frame}{Avancées depuis la réunion du 6 octobre}
  \begin{itemize}\espacement
    \item \textbf{État de l'art rédigé} sur les attaques contre les ML-NIDS.
    \item \textbf{Premier site développé} pour construire une topologie réseau par glisser-déposer, dans l'esprit de Packet Tracer.
  \end{itemize}
  \vspace{1.2em}
  \begin{block}{À faire ensuite}
    Choisir les modèles et les attaques, puis les intégrer au démonstrateur.
  \end{block}
  \note{La détection et les attaques sont du travail à venir. Aucun choix de couple modèle--attaque n'est arrêté.}
\end{frame}

\begin{frame}{État de l'art : pourquoi trois nouveaux papiers ?}
  \textbf{Au 6 octobre, l'empoisonnement reposait sur Venkatesan et Das :}
  \begin{itemize}
    \item l'objectif était seulement de \textbf{dégrader} le détecteur (plus d'erreurs) ;
    \item l'attaquant était \textbf{déjà à l'intérieur} (accès aux données, contrôleur SDN compromis) ;
    \item les attaques portaient sur des \textbf{tableaux de caractéristiques}, pas sur des paquets.
  \end{itemize}
  \vspace{.6em}
  \textbf{Trois manques, trois papiers ajoutés :}
  \begin{enumerate}\espacement
    \item Un attaquant \textbf{qui n'entre pas} chez la victime $\rightarrow$ \textbf{Alhussien}.
    \item Les \textbf{portes dérobées} : le modèle semble bon, mais laisse passer le trafic qui porte un motif secret $\rightarrow$ \textbf{PCAP-Backdoor} et \textbf{Ma}.
    \item Une attaque construite sur de \textbf{vrais paquets (PCAP)}, comme ceux capturés par notre prototype $\rightarrow$ \textbf{PCAP-Backdoor}.
  \end{enumerate}
  \note{Venkatesan et Das étaient déjà présentés le 6 octobre : ils ne font pas partie des nouveaux articles. Une porte dérobée reste une forme d'empoisonnement : elle demande aussi d'agir sur l'entraînement. Envoyer du trafic à un modèle figé ne suffit pas.}
\end{frame}

\begin{frame}{État de l'art : ce que chaque nouveau papier apporte}
  \footnotesize
  \renewcommand{\arraystretch}{1.2}
  \begin{tabularx}{\linewidth}{@{}>{\raggedright\arraybackslash}p{.15\linewidth}*{3}{>{\raggedright\arraybackslash}X}@{}}
    \toprule
    & \textbf{Qui est l'attaquant ?} & \textbf{Ce qu'il fait} & \textbf{Pourquoi on l'a retenu} \\
    \midrule
    \textbf{\href{https://doi.org/10.1109/NOMS56928.2023.10154453}{Alhussien}} (2023)
      & Il pollue une \textbf{base de données publique} que la victime utilise pour entraîner son détecteur.
      & Il modifie 5\,\% de cette base. La détection des attaques passe de 79\,\% à 47\,\%.
      & Scénario \textbf{réaliste} : pas besoin de pirater la victime. Défense simple : plus de données internes. \\
    \textbf{\href{https://doi.org/10.1145/3805033}{PCAP-Backdoor}} (2026)
      & Il peut ajouter du trafic dans les \textbf{captures d'entraînement}.
      & Il glisse un motif dans du trafic normal, \textbf{sans changer les étiquettes}. Ses attaques avec ce motif sont ensuite classées « bénin ».
      & Travaille sur des \textbf{fichiers PCAP}, comme notre capture : le papier le plus proche du démonstrateur. \\
    \textbf{\href{https://doi.org/10.1049/ise2/6664900}{Ma}} (2025)
      & Le \textbf{prestataire} à qui la victime confie l'entraînement.
      & Il marque 0,1\,\% des exemples avec un motif et les étiquette « bénin ». Plus de 97\,\% des attaques marquées passent.
      & Montre qu'un \textbf{test classique ne suffit pas} (moins de 3\,\% de perte) et propose une \textbf{défense}. \\
    \bottomrule
  \end{tabularx}
  \vfill
  \normalsize
  \textbf{À retenir :} une porte dérobée ne se voit pas aux tests normaux, seulement avec le motif.
  \note{Numérotation des fichiers dans docs/attaques_concretes_ML_NIDS : 07 Das, 08 Venkatesan, 09 Chathoth (PCAP-Backdoor), 10 Ma, 11 Alhussien. Das était déjà présent : ne pas le compter parmi les nouveaux articles. Sources : section Empoisonnement et références de etat_de_l_art_attaques_ml_nids.tex ; selection_attaques.md. Das change les étiquettes via un contrôleur SDN compromis ; Venkatesan insère des exemples contradictoires et vise notamment les faux positifs. Alhussien déplace le point d'entrée vers une source externe utilisée en apprentissage few-shot. PCAP-Backdoor agit sur les paquets avec des étiquettes propres et associe le déclencheur à la classe bénigne. Ma considère un prestataire d'entraînement malveillant, modifie les caractéristiques et les étiquettes d'une petite partie des exemples ; les tests sans déclencheur peuvent ne pas révéler la porte dérobée. Une backdoor reste une forme d'empoisonnement : les catégories ne sont pas disjointes. Aucun choix d'intégration n'est arrêté.}
\end{frame}

\begin{frame}{Le site : construire une topologie réseau}
  \textbf{Une interface web inspirée de Packet Tracer.}
  \vspace{1em}
  \begin{itemize}\espacement
    \item Ajouter et déplacer des équipements par glisser-déposer.
    \item Relier les équipements aux switches avec l'outil câble.
    \item Modifier les paramètres réseau dans le panneau de propriétés.
    \item Vérifier la cohérence de la topologie avant le déploiement.
  \end{itemize}
  \note{Vérification dans le code : Palette.jsx et TopologyCanvas.jsx pour le glisser-déposer ; useCableTool.js pour les liens ; DeviceForm.jsx pour les propriétés ; TopologyEditor.jsx et validators.js pour les contrôles. Ces fonctions sont présentes dans le dépôt consulté ; leur fonctionnement dans l'installation du groupe n'a pas été vérifié ici. Ne pas assimiler la présence d'un équipement IDS à une détection ML opérationnelle.}
\end{frame}

\begin{frame}{Exécution et observation : fonctions implémentées}
  \textbf{Le code du prototype comprend :}
  \vspace{1em}
  \begin{itemize}\espacement
    \item Le déploiement de la topologie sous forme de conteneurs et de réseaux Docker.
    \item Un terminal pour les PC et les serveurs démarrés.
    \item La capture des paquets sur une interface d'équipement reliée à un switch.
    \item L'affichage des paquets, des filtres et un export PCAP.
  \end{itemize}
  \vfill
  La détection ML et l'exécution des attaques restent à intégrer.
  \note{Fonctions suivies dans le code, de l'interface au backend : useDeployment.js vers POST /api/deployments et LabManager.deploy ; TerminalView.jsx vers TerminalSession ; CaptureView.jsx vers CaptureSession. Le téléchargement PCAP est construit dans le navigateur. Cette vérification du code ne constitue pas un test de bout en bout dans l'installation du groupe.}
\end{frame}

\begin{frame}[plain]
  \centering
  {\fontsize{36}{44}\selectfont\bfseries\color{projectblue} Démonstration du prototype\par}
\end{frame}

\begin{frame}{Intégration du ML-NIDS : travail restant}
  \begin{enumerate}\espacement
    \item \textbf{Choisir et reproduire une référence :} fixer le modèle, l'extracteur et les données de départ.
    \item \textbf{Relier la sonde au détecteur :} transformer le trafic capturé en caractéristiques et obtenir un score ou une décision.
    \item \textbf{Intégrer une attaque compatible :} préciser les accès nécessaires et les limites de modification du trafic.
    \item \textbf{Afficher la comparaison :} trafic de référence, trafic transformé et résultat de détection.
  \end{enumerate}
  \vfill
  \textbf{Premier objectif proposé :} une expérience reproductible de bout en bout sur une topologie simple.
  \note{Environ 45 secondes. L'image de la sonde est encore provisoire dans backend/config/devices.toml. Éviter de multiplier les modèles avant d'avoir une première chaîne complète.}
\end{frame}

\begin{frame}{Questions pour Grégory Blanc}
  \begin{enumerate}\espacement
    \item \textbf{Comparaison avant / après :} quelles informations afficher pour rendre l'effet de l'attaque facile à comprendre ?
    \item \textbf{Représentation :} privilégier une vue du réseau, une chronologie des échanges ou des graphiques de résultats ?
    \item \textbf{Déroulement :} montrer l'attaque en direct, étape par étape, ou présenter un bilan comparatif à la fin ?
    \item \textbf{Message pédagogique :} que doit retenir l'utilisateur après avoir vu une démonstration ?
  \end{enumerate}
\end{frame}

\end{document}
